\documentclass[12pt]{amsart}
%prepared in AMSLaTeX, under LaTeX2e
\addtolength{\oddsidemargin}{-1.1in} 
\addtolength{\evensidemargin}{-1.1in}
\addtolength{\topmargin}{-.9in}
\addtolength{\textwidth}{1.75in}
\addtolength{\textheight}{1.75in}

\renewcommand{\baselinestretch}{1.05}

\usepackage{verbatim} % for "comment" environment

\usepackage{palatino}

\usepackage[final]{graphicx}

\usepackage{tikz}
\usetikzlibrary{positioning}

\usepackage{bm,enumitem,xspace,fancyvrb}

\newtheorem*{thm}{Theorem}
\newtheorem*{defn}{Definition}
\newtheorem*{example}{Example}
\newtheorem*{problem}{Problem}
\newtheorem*{remark}{Remark}

\DefineVerbatimEnvironment{mVerb}{Verbatim}{numbersep=2mm,frame=lines,framerule=0.1mm,framesep=2mm,xleftmargin=4mm,fontsize=\footnotesize}

% macros
\usepackage{amssymb}
\newcommand{\bA}{\mathbf{A}}
\newcommand{\bB}{\mathbf{B}}
\newcommand{\bE}{\mathbf{E}}
\newcommand{\bF}{\mathbf{F}}
\newcommand{\bJ}{\mathbf{J}}
\newcommand{\bX}{\mathbf{X}}
\newcommand{\bL}{\mathbf{\Lambda}}

\newcommand{\bb}{\mathbf{b}}
\newcommand{\bc}{\mathbf{c}}
\newcommand{\bd}{\mathbf{d}}
\newcommand{\bi}{\mathbf{i}}
\newcommand{\bj}{\mathbf{j}}
\newcommand{\bk}{\mathbf{k}}
\newcommand{\br}{\mathbf{r}}
\newcommand{\bv}{\mathbf{v}}
\newcommand{\bw}{\mathbf{w}}
\newcommand{\bx}{\mathbf{x}}

\newcommand{\bzero}{\bm{0}}

\newcommand{\eps}{\epsilon}
\newcommand{\grad}{\nabla}
\newcommand{\ip}[2]{\ensuremath{\left<#1,#2\right>}}
\newcommand{\lam}{\lambda}
\newcommand{\lap}{\triangle}

\newcommand{\Null}{\operatorname{null}}
\newcommand{\rank}{\operatorname{rank}}
\newcommand{\range}{\operatorname{range}}
\newcommand{\trace}{\operatorname{tr}}

\newcommand{\RR}{\mathbf{R}}

\newcommand{\ZZ}{\mathbb{Z}}

\newcommand{\prob}[1]{\bigskip\noindent\textbf{#1.}\quad }
\newcommand{\exer}[2]{\prob{Exercise #2 on page #1}}

\newcommand{\pts}[1]{(\emph{#1 pts}) }
\newcommand{\epart}[1]{\bigskip\noindent\textbf{(#1)}\quad }
\newcommand{\ppart}[1]{\,\textbf{(#1)}\quad }

\newcommand{\Matlab}{\textsc{Matlab}\xspace}
\newcommand{\ds}{\displaystyle}

\def\vectwo#1#2{\begin{bmatrix}#1\\#2\end{bmatrix}}
\def\vecthree#1#2#3{\begin{bmatrix}#1\\#2\\#3\end{bmatrix}}
\def\vecfour#1#2#3#4{\begin{bmatrix}#1\\#2\\#3\\#4\end{bmatrix}}
\def\bbm{\begin{bmatrix}}
\def\ebm{\end{bmatrix}}


\begin{document}
\scriptsize \noindent Math 314 Linear Algebra \hfill Section 6.2 Diagonalizing a Matrix 
\normalsize\medskip


\thispagestyle{empty}
\begin{enumerate}
\item Find the eigenvalues $\lambda_1 \leq \lambda_2 \leq \lambda_3$ and associated eigenvectors, $\bx_1,\bx_2,\bx_3$ for the matrix $\bA$ below.  \\

$\bA =  \begin{bmatrix} 2&1&0\\ 0&3&0\\0&1&1 \end{bmatrix} $
\vfill
\newpage
\item Define the two matrices $\bX = [ \bx_1\: \bx_2\: \bx_3]$ and $\bL$ to be a 3 by 3 \emph{diagonal} matrix with $\lambda_1,\lambda_2,\lambda_3$ down the main diagonal \emph{in that order} and write them below.\\
\vfill
	\begin{enumerate}
	\item Is $\bX$ invertible or not? Why?
	\vfill
	\item Compute $\bA \bX.$
	\vfill
	\item Compute $\bX \bL.$
	\vfill
	\item Explain, without computation, that $\bA=\bX \bL \bX^{-1}$
	\vfill
	\item Use your answer from part (d), to find another way to compute $\bA^5.$ Are there advantages to the alternative approach?
	\vfill
	\end{enumerate}
\end{enumerate}
\end{document}

